\documentclass[10pt]{article}
\usepackage[T1]{fontenc}
\usepackage{amsmath,amssymb,amsfonts}

\newcommand{\Fclass}{\mathbf{F}}
\newcommand{\defeq}{\mathrel{\stackrel{\mathrm{def}}{=}}}
\newcommand{\lcmop}{\operatorname{lcm}}

\begin{document}

\section*{3. THE METHOD OF EXPONENT PAIRS}

\subsection*{3.1 INTRODUCTION}

In this chapter, we will develop Phillips' method of exponent pairs. We are
again striving for upper bounds for sums of the form
\begin{equation*}
 S=\sum_{n\in I}e(f(n)), \tag{3.1.1}
\end{equation*}
where $I=(a,b]$. However, we now assume that there is some $N>0$ such that
$I\subseteq[N,2N]$, and our estimates will be in terms of $N$ rather than
$|I|$. We will also assume that $f$ is some suitably well-defined function.
The exact conditions on $f$ are complicated (see Section 3.3), but the most
important feature is that $f'(x)$ is approximately $yx^{-s}$ for some $y>0$
and some $s>0$. The bounds for $S$ are expressed in terms of $N$ and
$L\defeq yN^{-s}$. (Note that $f'\approx L$.) If $L\leq1$, then a satisfactory
estimate can be obtained from Theorems 2.1 and 2.2, so we concentrate on the
case $L\geq1$. In this case, we seek an estimate of the form
\[
 S\ll L^kN^l.
\]
If we can prove such an estimate, we say that $(k,l)$ is an exponent pair.

Clearly, $(0,1)$ is an exponent pair. Phillips' theory consists of two
processes---$A$ and $B$---which produce new exponent pairs from old ones.

The $A$-process is based on Lemma 2.5, the Weyl--van der Corput inequality. We
shall see that if $(k,l)$ is an exponent pair, then so is
\[
 A(k,l)=\left(\frac{k}{2k+2},\frac{k+l+1}{2k+2}\right).
\]

The $B$-process is derived from the Poisson summation formula. From the
exponent pair $(k,l)$, the $B$-process produces the pair
\[
 B(k,l)=(l-1/2,k+1/2).
\]
Note that $B(0,1)=(1/2,1/2)$; this exponent pair gives the same estimate as
Theorem 2.9 in the case $q=0$. Moreover, $AB(0,1)=(1/6,2/3)$, which is
Theorem 2.9 in the case $q=1$. (The product $AB$ is to be interpreted as a
composition in the usual sense of analysis. Thus
$AB(0,1)=A(1/2,1/2)=(1/6,2/3)$.)

The $B$-process is involutory; that is, $B^2(k,l)=(k,l)$. Consequently, any
exponent pair derivable from Phillips' theory has the form
\begin{equation*}
 A^{q_1}BA^{q_2}B\cdots A^{q_k}B(0,1) \tag{3.1.2}
\end{equation*}
or
\begin{equation*}
 BA^{q_1}BA^{q_2}B\cdots A^{q_k}B(0,1), \tag{3.1.3}
\end{equation*}
where the $q_i$ are positive integers.

In Chapter 4, we will give several applications of exponent pairs. In these
applications, the ultimate result will be expressed as a function of $k$ and
$l$ of the form
\begin{equation*}
 \frac{ak+bl+c}{dk+el+f}, \tag{3.1.4}
\end{equation*}
where $a,\ldots,f$ are real numbers. In Chapter 5, we consider the question of
how to minimize (3.1.4), or equivalently, how to make an optimal choice of the
$q_i$ in (3.1.2) or (3.1.3).

This chapter is organized in the following way. In Section 3.2, we give
several lemmas on exponential integrals that will be used in the proof of the
$B$-process. In Section 3.3, we give heuristic arguments for the $A$- and
$B$-processes. Using these heuristics as a guide, we give the formal definition
of an exponent pair. Sections 3.4 and 3.5 contain the proofs of the $A$- and
$B$-processes, respectively.

\subsection*{3.2 LEMMAS ON EXPONENTIAL INTEGRALS}

In this section, we prove a number of estimates on exponential integrals.
However, the only result that will be used in the other sections of this
chapter is Lemma 3.6, which is the main tool necessary to prove the
$B$-process.

\medskip
\noindent\textbf{LEMMA 3.1.}
Assume that $f$ and $g$ are twice differentiable on $[a,b]$. Assume moreover
that $g/f'$ is monotonic and that $|f'(x)/g(x)|\geq\lambda$. Then
\[
 \int_a^b g(x)e(f(x))\,dx\ll\frac1\lambda.
\]

\noindent\textit{Proof.}
Integration by parts shows that the left-hand side is
\[
\begin{aligned}
 \frac{1}{2\pi i}\int_a^b\frac{g(x)}{f'(x)}\,d(e(f(x)))
 &=\left.\frac{g(x)e(f(x))}{2\pi i f'(x)}\right|_a^b
   -\frac{1}{2\pi i}\int_a^b e(f(x))
       \frac{d}{dx}\left(\frac{g(x)}{f'(x)}\right)\,dx \\
 &\ll \frac1\lambda+
   \int_a^b\left|\frac{d}{dx}\left(\frac{g(x)}{f'(x)}\right)\right|\,dx.
\end{aligned}
\]
Since $g/f'$ is monotonic, the absolute-value bars may be taken outside the
last integral, and the desired estimate follows.

\medskip
\noindent\textbf{LEMMA 3.2.}\par\noindent
Assume that $f$ has two continuous derivatives on $[a,b]$ and that
$|f''(x)|\geq\lambda_2>0$ there. Then
\[
 \int_a^b e(f(x))\,dx\ll\lambda_2^{-1/2}.
\]

\noindent\textit{Proof.}
Let $\lambda$ be a number to be chosen later. We write
$[a,b]=E_1\cup E_2$, where
\[
 E_1=\{x:|f'(x)|\geq\lambda\}
 \quad\text{and}\quad
 E_2=\{x:|f'(x)|<\lambda\}.
\]
Since $f''$ has constant sign, $E_1$ consists of at most two intervals and
$E_2$ consists of at most one interval. By Lemma 3.1,
\[
 \int_{E_1}e(f(x))\,dx\ll\frac1\lambda.
\]
Let $E_2=(c,d)$. Then
\[
 (d-c)\lambda_2
 \leq\left|\int_c^d f''(x)\,dx\right|
 =|f'(d)-f'(c)|\leq2\lambda,
\]
so
\[
 \left|\int_{E_2}e(f(x))\,dx\right|\leq d-c\ll\frac{\lambda}{\lambda_2}.
\]
Combining the above gives
\[
 \int_a^b e(f(x))\,dx\ll\frac1\lambda+\frac{\lambda}{\lambda_2}.
\]
Choosing $\lambda=\lambda_2^{1/2}$ completes the proof.

\medskip
\noindent\textbf{LEMMA 3.3.}
Let $A$ and $X$ be positive numbers. Then
\[
 \int_{-X}^{X}e(Ax^2)\,dx
 =\frac{e(1/8)}{\sqrt{2A}}+O\left(\frac1{AX}\right).
\]

\noindent\textit{Proof.}
After a couple of obvious changes of variables, we see that it suffices to
show that
\[
 \int_0^X e(x^2)\,dx
 =\frac{e(1/8)}{2\sqrt2}+O\left(\frac1X\right).
\]
Consider the contour integral
\[
 \int_{\mathbf C}e(z^2)\,dz,
\]
where $\mathbf C$ is the path that goes along the straight line from $0$ to
$X$, along the circular arc from $X$ to $Xe(1/8)$, and along the straight line
from $Xe(1/8)$ to $0$. Since the integrand is entire, the integral around the
contour is zero. Therefore,
\[
 \int_0^X e(x^2)\,dx+\int_{\mathbf C_1}e(z^2)\,dz
 =e(1/8)\int_0^X e^{-2\pi x^2}\,dx,
\]
where $\mathbf C_1$ is the circular arc of $\mathbf C$. Since
$\sin 2\theta\geq4\theta/\pi$ for $0\leq\theta\leq\pi/4$, we have
\[
\begin{aligned}
 \left|\int_{\mathbf C_1}e(z^2)\,dz\right|
 &\leq X\int_0^{\pi/4}\exp(-2\pi X^2\sin 2\theta)\,d\theta \\
 &\leq X\int_0^{\pi/4}\exp(-8X^2\theta)\,d\theta
 \ll\frac1X.
\end{aligned}
\]
Furthermore,
\[
\begin{aligned}
 \int_0^X e^{-2\pi x^2}\,dx
 &=\frac{1}{2\sqrt{2\pi}}\int_0^{2\pi X^2}y^{-1/2}e^{-y}\,dy \\
 &=\frac{1}{2\sqrt{2\pi}}
   \left\{\Gamma(1/2)-\int_{2\pi X^2}^{\infty}y^{-1/2}e^{-y}\,dy\right\} \\
 &=\frac{1}{2\sqrt2}+O\left(\frac1X\right).
\end{aligned}
\]
The result follows.

\medskip
\noindent\textbf{LEMMA 3.4.}
Suppose $g$ is a real-valued function with four continuous derivatives on
$[a,b]$. Suppose also that $g''(x)\geq\lambda_2>0$ and that $g'(x_0)=0$ for
some $x_0\in[a,b]$. Finally, assume that there are positive constants
$\lambda_3$ and $\lambda_4$ such that
\[
 |g^{(3)}(x)|\leq\lambda_3
 \quad\text{and}\quad
 |g^{(4)}(x)|\leq\lambda_4
\]
on $[a,b]$. Then
\begin{equation*}
 \int_a^b e(g(x))\,dx
 =\frac{e(1/8+g(x_0))}{g''(x_0)^{1/2}}+O(R_1+R_2), \tag{3.2.1}
\end{equation*}
where
\[
 R_1=
 \min\left(\frac{1}{\lambda_2(x_0-a)},\frac1{\lambda_2^{1/2}}\right)
 +\min\left(\frac{1}{\lambda_2(b-x_0)},\frac1{\lambda_2^{1/2}}\right)
\]
and
\[
 R_2=(b-a)\lambda_4\lambda_2^{-2}
     +(b-a)\lambda_3^2\lambda_2^{-3}.
\]
If the condition $g''(x)\geq\lambda_2>0$ is replaced by
$g''(x)\leq-\lambda_2<0$, then we may replace $g$ by $-g$ and get the same
result, save that the main term is now
\[
 \frac{e(-1/8+g(x_0))}{|g''(x_0)|^{1/2}}.
\]
In our applications of Lemma 3.4, we will be assuming that
$\lambda_j\approx FN^{-j}$ for $j=2,3,4$. In this case, $R_2=O(F^{-1}N)$.

\noindent\textit{Proof.}
We may assume that
\begin{equation*}
 a+\lambda_2^{-1/2}\leq x_0\leq b-\lambda_2^{-1/2}, \tag{3.2.2}
\end{equation*}
for otherwise Lemma 3.2 is better than our claimed result.

We multiply both sides of (3.2.1) by $e(-g(x_0))$ to reduce to the case
$g(x_0)=0$. Similarly, we may translate coordinates so that $x_0=0$. Let
$q(x)=\tfrac12g''(0)x^2$ and $r(x)=g(x)-q(x)$. Since $q(x)$ is the Taylor
approximation to $g$ of order $2$, we may write
\begin{equation*}
 r(x)=\frac12x^3\int_0^1(1-u)^2g^{(3)}(xu)\,du. \tag{3.2.3}
\end{equation*}
Moreover, Taylor's theorem applied to $g'$ gives
\begin{equation*}
 r'(x)=x^2\int_0^1(1-u)g^{(3)}(xu)\,du. \tag{3.2.4}
\end{equation*}
We now write
\[
 \int_a^b e(g(x))\,dx
 =\int_a^b e(q(x))\,dx
  +\int_a^b e(q(x))\{e(r(x))-1\}\,dx.
\]
The first integral can be handled with Lemma 3.3. It then suffices to show
that
\begin{equation*}
 \int_a^b e(q(x))\{e(r(x))-1\}\,dx\ll R_1+R_2. \tag{3.2.5}
\end{equation*}
Integrating by parts, we see that this integral is
\begin{align*}
 &\left.\frac{e(q(x))\{e(r(x))-1\}}{2\pi i g''(0)x}\right|_a^b \tag{3.2.6}\\
 &\qquad
 -\frac1{g''(0)}\int_a^b e(q(x))
   \left\{\frac{e(r(x))r'(x)}x
   -\frac{e(r(x))-1}{2\pi i x^2}\right\}\,dx.
\end{align*}
The first term contributes
$\ll |a|^{-1}\lambda_2^{-1}+b^{-1}\lambda_2^{-1}\ll R_1$. We write the
% ed.: the printed text has ``as $T_1+\frac1{2\pi}T_2$''; with $T_1=\int e(g)r'/x$
% ed.: and $T_2=\int e(q)(e(r)-1)/x^2$ as used below, the integral in (3.2.6) is
% ed.: $-g''(0)^{-1}\{T_1-(2\pi i)^{-1}T_2\}$, which is what (3.2.7) uses.
integral in (3.2.6) as $-\frac1{g''(0)}\left\{T_1-\frac1{2\pi i}T_2\right\}$,
where $T_1$ arises from the first term in braces and $T_2$ from the second.

Let $h(x)=r'(x)x^{-2}$ and $j(x)=g'(x)x^{-1}$. Then
\[
\begin{aligned}
 T_1
 &=\int_a^b\frac{e(g(x))r'(x)}x\,dx
  =\int_a^b\frac{h(x)}{j(x)}e(g(x))g'(x)\,dx \\
 &=\frac1{2\pi i}\int_a^b\frac{h(x)}{j(x)}\,d(e(g(x))) \\
 &=\left.\frac{h(x)e(g(x))}{2\pi i j(x)}\right|_a^b
   -\frac1{2\pi i}\int_a^b
       \frac{d}{dx}\left(\frac{h(x)}{j(x)}\right)e(g(x))\,dx.
\end{aligned}
\]
From (3.2.4), we see that
\[
 h(x)=\int_0^1(1-u)g^{(3)}(xu)\,du\ll\lambda_3
\]
and
\[
 h'(x)=\int_0^1(1-u)u g^{(4)}(xu)\,du\ll\lambda_4.
\]
Similarly,
\[
 j(x)=\int_0^1g''(xu)\,du\geq\lambda_2,
 \qquad
 j'(x)=\int_0^1u g^{(3)}(xu)\,du\ll\lambda_3.
\]
Thus
\[
 \frac{d}{dx}\left(\frac hj\right)
 =\frac{h'}j-\frac{hj'}{j^2}
 \ll\frac{\lambda_4}{\lambda_2}
     +\frac{\lambda_3^2}{\lambda_2^2},
\]
and so
\[
 T_1\ll\frac{\lambda_3}{\lambda_2}
 +(b-a)\frac{\lambda_4}{\lambda_2}
 +(b-a)\frac{\lambda_3^2}{\lambda_2^2}.
\]

To treat $T_2$, let $\delta$ be a parameter to be chosen later, and set
\[
 \mathbf I=[a,b]\cap[-\delta,\delta],
 \qquad
 \mathbf J=[a,b]\setminus[-\delta,\delta].
\]
Note that $\mathbf I=[c,d]$, where $c=\max(a,-\delta)$ and
$d=\min(b,\delta)$. Define $k(x)=r(x)x^{-3}$. From (3.2.3), we see that
$k(x)\ll\lambda_3$ and $k'(x)\ll\lambda_4$. Set
$m(x)=(e(x)-1)/x$; then $m(x)\ll1$ and $m'(x)\ll1$. The contribution of the
interval $\mathbf I$ to $T_2$ is
\[
\begin{aligned}
 \int_c^d e(q(x))x m(r(x))k(x)\,dx
 &=\left.\frac{e(q(x))m(r(x))k(x)}{2\pi i g''(0)}\right|_c^d \\
 &\quad-\frac1{2\pi i g''(0)}\int_c^d e(q(x))m'(r(x))r'(x)k(x)\,dx \\
 &\quad-\frac1{2\pi i g''(0)}\int_c^d e(q(x))m(r(x))k'(x)\,dx \\
 &\ll \lambda_2^{-1}\lambda_3
 +\lambda_2^{-1}\lambda_3^2\delta^3
 +\lambda_2^{-1}\lambda_4(b-a).
\end{aligned}
\]
For the interval $\mathbf J$, we apply Lemma 3.1 to the integrals
\[
 \int_{\mathbf J}\frac{e(g(x))}{x^2}\,dx
 \quad\text{and}\quad
 \int_{\mathbf J}\frac{e(q(x))}{x^2}\,dx.
\]
In the first integral, we are led to consider the function $g'(x)x^2$. Its
absolute value is at least $\lambda_2\delta^3$, and it is monotonic since
\[
 (g'(x)x^2)'=g''(x)x^2+2xg'(x)>0.
\]
By Lemma 3.1,
\[
 \int_{\mathbf J}\frac{e(g(x))}{x^2}\,dx
 \ll\lambda_2^{-1}\delta^{-3}.
\]
Similarly, since $q'(x)x^2=g''(0)x^3$ is monotonic,
\[
 \int_{\mathbf J}\frac{e(q(x))}{x^2}\,dx
 \ll\lambda_2^{-1}\delta^{-3}.
\]
We take $\delta=\lambda_3^{-1/3}$ and combine this with our estimate of
$T_1$ to see that the integral in (3.2.5) is
\begin{equation*}
 \begin{aligned}
 &\ll R_1+\lambda_2^{-1}(T_1+T_2) \\
 &\ll R_1+\lambda_2^{-2}\lambda_3
 +(b-a)\lambda_2^{-2}\lambda_4
 +(b-a)\lambda_2^{-3}\lambda_3^2.
 \end{aligned} \tag{3.2.7}
\end{equation*}
Put $U=\lambda_2^{-1}(b-a)^{-1}$. From (3.2.2), we see that $U\ll R_1$.
Since
\[
 \lambda_2^{-2}\lambda_3
 =U^{1/2}\{\lambda_2^{-3}\lambda_3^2(b-a)\}^{1/2}
 \ll\max\{R_1,\lambda_2^{-3}\lambda_3^2(b-a)\},
\]
the second term in (3.2.7) may be omitted. This completes the proof.

For the next lemma, recall that we defined $[x]$ to be the greatest integer
at most $x$, and $\psi(x)=x-[x]-1/2$.

\medskip
\noindent\textbf{LEMMA 3.5.}
Suppose $f$ has two continuous derivatives on $[a,b]$. Suppose also that $f'$
is decreasing, that $H_1$ and $H_2$ are integers such that
$H_1<f'(x)<H_2$, and that $H=H_2-H_1\geq2$. Then
\[
 \sum_{n\in I}e(f(n))
 =\sum_{H_1\leq h\leq H_2}\int_a^b e(f(x)-hx)\,dx+O(\log H).
\]

\noindent\textit{Proof.}
The left-hand side is unchanged if we replace $f(n)$ by $f(n)-kn$ for any
% ed.: the original normalizes only to $H_1<0<H_2$; the bound $\ll H_2/(h-H_2)$
% ed.: below uses $|f'|<H_2$ (and symmetrically $|f'|<-H_1$), which requires the
% ed.: balanced choice $k=[(H_1+H_2)/2]$, giving $-H_1\asymp H_2\asymp H$.
integer $k$. Therefore, taking $k=[(H_1+H_2)/2]$, we may normalize to the
case $H_1<0<H_2$ with $-H_1\asymp H_2\asymp H$.

Using Riemann--Stieltjes integration, we get
\begin{equation*}
 \sum_{n\in I}e(f(n))
 =\int_a^b e(f(x))\,dx-\int_a^b e(f(x))\,d\psi(x)+O(1). \tag{3.2.8}
\end{equation*}
Integration by parts shows that
\[
 -\int_a^b e(f(x))\,d\psi(x)
 =2\pi i\int_a^b e(f(x))f'(x)\psi(x)\,dx+O(1).
\]
The symmetric partial sums of the Fourier expansion
\[
 \psi(x)=\sum_{h\neq0}\frac{e(-hx)}{2\pi i h}
\]
are uniformly bounded and converge to $\psi(x)$ away from the integers; the
values at the integers do not affect the integral. Hence
\begin{equation*}
 2\pi i\int_a^b e(f(x))f'(x)\psi(x)\,dx
 =\sum_{h\neq0}\frac1h\int_a^b e(f(x)-hx)f'(x)\,dx. \tag{3.2.9}
\end{equation*}
If $h<H_1$ or $h>H_2$, then $f'(x)-h$ is monotonic and nonzero on $[a,b]$.
Moreover, $f'(x)/(f'(x)-h)$ is monotonic, so Lemma 3.1 gives
\[
 \int_a^b e(f(x)-hx)f'(x)\,dx
 \ll\frac{|f'(b)|}{|f'(b)-h|}+\frac{|f'(a)|}{|f'(a)-h|}.
\]
When $h>H_2$, the above is $\ll H_2/(h-H_2)$. Thus these terms contribute
\[
 \ll\sum_{\substack{h>H_2\\h\neq0}}\frac{H_2}{h(h-H_2)}
 \ll\sum_{\substack{h>H_2\\h\neq0}}
 \left\{\frac1{h-H_2}-\frac1h\right\}
 \ll\log H
\]
to (3.2.9). The same upper bound holds for the terms with $h<H_1$. We
integrate the remaining terms by parts to get
\[
\begin{aligned}
 &\sum_{\substack{H_1\leq h\leq H_2\\h\neq0}}
   \frac1h\int_a^b f'(x)e(f(x)-hx)\,dx \\
 &=\sum_{\substack{H_1\leq h\leq H_2\\h\neq0}}
   \frac1{2\pi i h}\int_a^b e(-hx)\,d(e(f(x))) \\
 &=\left.\sum_{\substack{H_1\leq h\leq H_2\\h\neq0}}
   \frac{e(f(x)-hx)}{2\pi i h}\right|_a^b
   +\sum_{\substack{H_1\leq h\leq H_2\\h\neq0}}
   \int_a^b e(f(x)-hx)\,dx \\
 &=\sum_{\substack{H_1\leq h\leq H_2\\h\neq0}}
   \int_a^b e(f(x)-hx)\,dx+O(\log H).
\end{aligned}
\]
Combining the above with (3.2.8) and (3.2.9) yields
\[
 \sum_{n\in I}e(f(n))
 =\int_a^b e(f(x))\,dx
 +\sum_{\substack{H_1\leq h\leq H_2\\h\neq0}}
  \int_a^b e(f(x)-hx)\,dx+O(\log H).
\]
Since $H_1<0<H_2$, this completes the proof.

\medskip
\noindent\textbf{LEMMA 3.6.}
Suppose that $f$ has four continuous derivatives on $[a,b]$, and that
$f''<0$ on this interval. Suppose further that $[a,b]\subseteq[N,2N]$ and
that $\alpha=f'(b)$ and $\beta=f'(a)$. Assume that there is some $F>0$ such
that
\[
 -f''(x)\approx FN^{-2},
 \qquad f^{(3)}(x)\ll FN^{-3},
 \qquad f^{(4)}(x)\ll FN^{-4}
\]
for $x\in[a,b]$. For $\nu\in[\alpha,\beta]$, let $x_\nu$ be defined by the
relation $f'(x_\nu)=\nu$, and let
$\phi(\nu)=-f(x_\nu)+\nu x_\nu$. Then
\[
 \sum_{n\in I}e(f(n))
 =\sum_{\alpha\leq\nu\leq\beta}
  \frac{e(-\phi(\nu)-1/8)}{|f''(x_\nu)|^{1/2}}
 +O\left(\log(FN^{-1}+2)+F^{-1/2}N\right).
\]

\noindent\textit{Proof.}
We may assume that $F\geq1$, for the theorem is trivial otherwise. Let
$H_1=[\alpha]-1$ and $H_2=[\beta]+1$. From the previous two lemmas, we see
that
\begin{equation*}
 \sum_{n\in I}e(f(n))
 =\sum_{\alpha<\nu<\beta}
  \frac{e(-\phi(\nu)-1/8)}{|f''(x_\nu)|^{1/2}}+O(E), \tag{3.2.10}
\end{equation*}
where
% ed.: the printed $E$ omits the $O(\log H)$ of Lemma 3.5, $H=H_2-H_1\ll FN^{-1}+1$;
% ed.: it is $\ll\log(FN^{-1}+2)$ and is absorbed in the final bound for $E$.
\[
\begin{aligned}
 E=\log(H_2-H_1)+\sum_{H_1\leq\nu\leq H_2}\bigl\{&F^{-1}N
 +\min\bigl(|f'(a)-\nu|^{-1},F^{-1/2}N\bigr)\\
 &+\min\bigl(|f'(b)-\nu|^{-1},F^{-1/2}N\bigr)\bigr\}.
\end{aligned}
\]
Note that
\[
 H_2-H_1\ll\left|\int_a^b f''(x)\,dx\right|+1\ll FN^{-1}+1,
\]
so
\[
 \sum_{H_1\leq\nu\leq H_2}F^{-1}N
 \ll F^{-1}N(FN^{-1}+1)\ll1+F^{-1}N.
\]
The non-harmonic remainder here is majorized by $F^{-1/2}N$ (and the
constant term is absorbed by the logarithm). Note also that, for
$\nu<H_2-1$,
\[
 |f'(a)-\nu|\geq H_2-1-\nu,
\]
so
\[
\begin{aligned}
 &\sum_{H_1\leq\nu\leq H_2}
 \min\bigl(|f'(a)-\nu|^{-1},F^{-1/2}N\bigr) \\
 &\qquad\ll F^{-1/2}N
 +\sum_{H_1\leq\nu<H_2-1}(H_2-1-\nu)^{-1} \\
 &\qquad\ll F^{-1/2}N+\log(FN^{-1}+2).
\end{aligned}
\]
The same upper bound holds for the sum involving $f'(b)$. Thus
\[
 E\ll\log(FN^{-1}+2)+F^{-1/2}N.
\]
Finally, extending the main sum in (3.2.10) to
$\alpha\leq\nu\leq\beta$ changes it by at most two terms, each of size
$O(F^{-1/2}N)$, and these may be absorbed into the error term.

\subsection*{3.3 HEURISTIC ARGUMENTS AND DEFINITIONS}

We begin this section by giving some heuristic arguments for the $A$- and
$B$-processes. Then, partly guided by these heuristics, we give the formal
definition of exponent pairs.

Let $S$ be as defined in (3.1.1). Recall the Weyl--van der Corput inequality,
which states that if $H\leq N$, then
\[
 |S|^2\ll\frac{N^2}{H}+\frac NH\sum_{1\leq h\leq H}|S_1(h)|,
\]
where
\begin{equation*}
 S_1(h)=\sum_{\substack{a<n\leq b\\a<n+h\leq b}}e(f_1(n;h)), \tag{3.3.1}
\end{equation*}
and $f_1(n;h)=f(n)-f(n+h)$. If $f'\approx L$, then
$f_1'\approx L|h|N^{-1}$. If we assume that the exponent pair $(k,l)$ can be
applied to $S_1(h)$, then we obtain
\[
\begin{aligned}
 |S|^2
 &\ll N^2H^{-1}+NH^{-1}\sum_{1\leq h\leq H}
      (L|h|N^{-1})^kN^l \\
 &\ll N^2H^{-1}+H^kL^kN^{1-k+l}.
\end{aligned}
\]
Upon choosing $H$ to make the two terms equal, we get
\[
 |S|\ll L^{k/(2k+2)}N^{(k+l+1)/(2k+2)}.
\]

For the $B$-process, we use Lemma 3.6 to get
\begin{equation*}
 S=e(-1/8)\sum_{\alpha\leq\nu\leq\beta}
 \frac{e(-\phi(\nu))}{|f''(x_\nu)|^{1/2}}+\text{error terms}. \tag{3.3.2}
\end{equation*}
For the purposes of this heuristic argument, we will ignore the error terms.
Now $\phi'(\nu)=x_\nu\approx N$, and
$|f''(x_\nu)|\approx LN^{-1}$. By partial summation (and complex conjugation),
we see that
\[
 \left|\sum_{\alpha\leq\nu\leq\beta}
 \frac{e(-\phi(\nu))}{|f''(x_\nu)|^{1/2}}\right|
 \ll L^{-1/2}N^{1/2}
 \max_{\alpha\leq\beta'\leq\beta}
 \left|\sum_{\alpha<\nu\leq\beta'}e(\phi(\nu))\right|.
\]
If we assume that the exponent pair $(k,l)$ may be applied to the last sum,
we get
\[
 |S|\ll L^{-1/2}N^{1/2}N^kL^l
 =L^{l-1/2}N^{k+1/2}.
\]

To make these arguments rigorous, we need to work with a suitable class of
functions $\Fclass$. This class should have the property that if
$f\in\Fclass$, then the auxiliary functions $f_1$ and $\phi$ of (3.3.1) and
(3.3.2), respectively, are also in $\Fclass$.

Another requirement for our class of functions is that it contain those
functions $f$ which satisfy $f'(x)=yx^{-s}$ for some $s>0$. For these
functions we have
\[
 f^{(p+1)}(x)=(-1)^p(s)_p yx^{-s-p}.
\]
Here, $(s)_p$ is defined recursively by the relations
\[
 (s)_0=1,
 \qquad
 (s)_{p+1}=(s+p)(s)_p.
\]
Note that when $p\geq1$, $(s)_p=s(s+1)\cdots(s+p-1)$.

\medskip
\noindent\textbf{Definition.}
Let $N,y,s$, and $\epsilon$ be positive real numbers with $\epsilon<1/2$, and
let $P$ be a nonnegative integer. Define $\Fclass(N,P,s,y,\epsilon)$ to be the
set of functions $f$ such that
\begin{enumerate}
 \item $f$ is defined and has $P$ continuous derivatives on some interval
       $[a,b]$, with $[a,b]\subseteq[N,2N]$;
 \item if $0\leq p\leq P-1$ and $a\leq x\leq b$, then
 \begin{equation*}
  \left|f^{(p+1)}(x)-(-1)^p(s)_p yx^{-s-p}\right|
  <\epsilon(s)_p yx^{-s-p}. \tag{3.3.3}
 \end{equation*}
\end{enumerate}

\medskip
\noindent\textbf{Definition.}
Let $k$ and $l$ be real numbers such that
$0\leq k\leq1/2\leq l\leq1$. Suppose that for every $s>0$, there is some
$P=P(k,l,s)$ and some $\epsilon=\epsilon(k,l,s)<1/2$ such that for every
$N>0$, every $y>0$, and every
$f\in\Fclass(N,P,s,y,\epsilon)$, the estimate
\begin{equation*}
 \left|\sum_{n\in I}e(f(n))\right|
 \ll (yN^{-s})^kN^l+y^{-1}N^s \tag{3.3.4}
\end{equation*}
holds. Here it is also assumed that $f$ is defined on $[a,b]$, and the implied
constant depends only on $k,l$, and $s$. We then say that $(k,l)$ is an
exponent pair.

In proving that $(k,l)$ is an exponent pair, we may always assume that
% ed.: Theorems 2.1 and 2.2 need $f'$ monotonic, resp.\ $|f''|\asymp yN^{-s-1}$;
% ed.: both follow from (3.3.3) with $p=1$, so $P\geq2$ is tacitly assumed here
% ed.: (for $P=1$ the class $\Fclass$ imposes nothing on $f''$ and the bounds fail).
$yN^{-s}\geq1$. For if $yN^{-s}<1/2$, then Theorem 2.1 (applicable since
we may take $P\geq2$, so that $f'$ is monotonic) yields
\[
 \left|\sum_{n\in I}e(f(n))\right|\ll y^{-1}N^s.
\]
If $1/2\leq yN^{-s}<1$, then Theorem 2.2 yields
\[
 \left|\sum_{n\in I}e(f(n))\right|\ll N^{1/2},
\]
and $N^{1/2}\ll(yN^{-s})^kN^l$ since $l\geq1/2$.

We close this section with an explanation of why the conditions
$0\leq k\leq1/2\leq l\leq1$ are included in the definition of exponent
pairs. The condition $l\geq1/2$ is no essential restriction, since no pair
$(k,l)$ which satisfies (3.3.4) can have $l<1/2$. To see why, consider, for a positive integer $N$, the sum
\[
 S(t)=\sum_{N<n\leq2N}e\left(\frac tn\right).
\]
By Theorem 2 of Appendix A,
\[
 \int_T^{2T}|S(t)|^2\,dt\geq(T-4N^2)N.
\]
It follows that if $T\geq8N^2$, then there is some $t$ with
$T\leq t\leq2T$ such that $|S(t)|\gg N^{1/2}$.

We also claim that there is a sequence $\{t_\nu\}_{\nu=1}^{\infty}$ such that
$t_\nu\to\infty$ and $|S(t_\nu)|=N$. Take
\[
 t_\nu=\nu\,\lcmop\{1,2,\ldots,2N\}.
\]
Then $t_\nu/n\in\mathbb Z$ whenever $N<n\leq2N$, and so $S(t_\nu)=N$.
This shows that no exponent pair $(k,l)$ can have $k<0$, and any exponent pair
that has $k=0$ must have $l=1$.

There is no need ever to take $l>1$ in an exponent pair, for an exponent pair
with $k\geq0$ and $l>1$ would yield an estimate worse than the trivial
estimate provided by the exponent pair $(0,1)$. Similarly, there is no need
to take $k>1/2$, for then $(k,l)$ would provide a weaker estimate than the
pair $(1/2,1/2)$ of Theorem 2.2.

Finally, we observe that if $(k,1/2)$ is an exponent pair, then $k=1/2$.
Assume that $(k,1/2)$ is an exponent pair. Let $H$ be a positive integer. Let
$t$ be defined by the equation
\[
 t^2=\lcmop\{1,2,\ldots,H\},
\]
and let $N=t^2H^{-2}$. Define
\[
 f(x)=2t x^{1/2}.
\]
Then $tN^{-1/2}=H$, and (3.3.4) gives
\begin{equation*}
 \left|\sum_{N<n\leq2N}e(f(n))\right|
 \ll H^kN^{1/2}+H^{-1}. \tag{3.3.5}
\end{equation*}
On the other hand, Lemma 3.6 gives
\[
\begin{aligned}
 \sum_{N<n\leq2N}e(f(n))
 &=\sum_{2^{-1/2}tN^{-1/2}<\nu\leq tN^{-1/2}}
  \frac{e(t^2/\nu-1/8)}{|f''(x_\nu)|^{1/2}}\\
 &\quad+O\left(\log(H+2)+H^{-1/2}N^{1/2}\right).
\end{aligned}
\]
Note that $tN^{-1/2}=H$ and $|f''(x_\nu)|\approx H/N$. Our definition of
$t$ ensures that $t^2/\nu\in\mathbb Z$ for every integer $\nu$ in the range
of summation. Therefore, for sufficiently large $H$,
\begin{equation*}
 \left|\sum_{N<n\leq2N}e(f(n))\right|
 \gg H^{1/2}N^{1/2}. \tag{3.3.6}
\end{equation*}
Estimates (3.3.5) and (3.3.6) are inconsistent unless $k\geq1/2$, so the
result is proved.

\subsection*{3.4 PROOF OF THE A-PROCESS}

Our first lemma shows that the auxiliary function $f_1$ of (3.3.1) is in an
appropriate $\Fclass$-set.

\medskip
\noindent\textbf{LEMMA 3.7.}
Suppose $P\geq1$, $f\in\Fclass(N,P,s,y,\epsilon)$, and $f$ is defined on
$[a,b]$. Suppose also that
\[
 1\leq h<\min\left(b-a,\frac{2\epsilon N}{s+P}\right).
\]
Define
\[
 f_1(x)=f(x)-f(x+h).
\]
Then $f_1$ is defined on $[a,b-h]$ and
% ed.: the Definition of \S3.3 requires $\epsilon<1/2$; for $3\epsilon\geq1/2$ the
% ed.: conclusion is to be read as (3.3.3) with $3\epsilon$ in place of $\epsilon$
% ed.: (Theorem 3.8 only uses $\epsilon_1=\epsilon/3$, so no issue arises there).
\[
 f_1\in\Fclass(N,P-1,s+1,shy,3\epsilon).
\]

\noindent\textit{Proof.}
Let
\[
 F(x)=
 \begin{cases}
  yx^{1-s}(1-s)^{-1},&s\neq1,\\
  y\log x,&s=1.
 \end{cases}
\]
Condition (3.3.3) may be written as
\[
 |f^{(p+1)}(x)-F^{(p+1)}(x)|
 <\epsilon|F^{(p+1)}(x)|.
\]
Define $F_1(x)=F(x)-F(x+h)$ and $G(x)=-hyx^{-s}$. We will show that
\begin{equation*}
 |f_1^{(p+1)}(x)-F_1^{(p+1)}(x)|
 <\epsilon|F_1^{(p+1)}(x)| \tag{3.4.1}
\end{equation*}
and that
\begin{equation*}
 |F_1^{(p+1)}(x)-G^{(p+1)}(x)|
 <\epsilon|G^{(p+1)}(x)| \tag{3.4.2}
\end{equation*}
when $0\leq p\leq P-2$. The desired result will then follow readily.

From our definitions of $f_1$ and $F_1$, we see that
\[
 |f_1^{(p+1)}(x)-F_1^{(p+1)}(x)|
 =\left|\int_x^{x+h}
   \{f^{(p+2)}(u)-F^{(p+2)}(u)\}\,du\right|.
\]
When $a\leq x\leq b-h$ and $0\leq p\leq P-2$, the above is bounded by
\[
 \epsilon\int_x^{x+h}|F^{(p+2)}(u)|\,du
 =\epsilon|F_1^{(p+1)}(x)|.
\]
This proves (3.4.1). For (3.4.2), we observe that the left-hand side is equal
to
\[
 y(s)_{p+1}\left|\int_x^{x+h}
  (u^{-s-p-1}-x^{-s-p-1})\,du\right|
 =y(s)_{p+2}\left|\int_x^{x+h}\int_x^u
  w^{-s-p-2}\,dw\,du\right|,
\]
and this in turn is
\[
 \leq\frac12(s)_{p+2}yh^2x^{-s-p-2}
 <\epsilon(s+1)_p shy x^{-s-p-1},
\]
provided $h<2\epsilon x/(s+p+1)$. This last condition is guaranteed by our
hypothesis $h<2\epsilon N/(s+P)$.

To complete the proof, we note that
\[
 |F_1^{(p+1)}(x)|<(1+\epsilon)|G^{(p+1)}(x)|
\]
follows from (3.4.2). Combining this with (3.4.1) and (3.4.2) yields
\[
 |f_1^{(p+1)}(x)-G^{(p+1)}(x)|
 <(2\epsilon+\epsilon^2)|G^{(p+1)}(x)|.
\]
Since $2\epsilon+\epsilon^2<3\epsilon$ for $\epsilon<1$, the proof is
complete.

\medskip
\noindent\textbf{THEOREM 3.8.}
If $(k,l)$ is an exponent pair, then
\[
 (\kappa,\lambda)=A(k,l)
 =\left(\frac{k}{2k+2},\frac{k+l+1}{2k+2}\right)
\]
is an exponent pair.

\noindent\textit{Proof.}
First, since $0\leq k$ and $1/2\leq l\leq1$, we have
\[
 0\leq\kappa=\frac{k}{2k+2}<\frac{k+1}{2k+2}=\frac12,
 \qquad
 \frac12\leq\lambda=\frac12+\frac{l}{2k+2}\leq1.
\]
Now let $y,N$, and $s$ be positive. We need to show that there exist
$P_1>0$ and $\epsilon_1>0$ such that if
$f\in\Fclass(N,P_1,s,y,\epsilon_1)$ and $L=yN^{-s}\geq1$, then
\[
 \left|\sum_{n\in I}e(f(n))\right|\ll L^\kappa N^\lambda.
\]
Our proof breaks into two cases: (i) $L\geq\log N$ and
(ii) $1\leq L<\log N$. We begin with case (i).

Since $(k,l)$ is an exponent pair, for the parameter $s+1$ there exist
$P>0$ and $\epsilon$ with $0<\epsilon<1/2$ such that the defining exponent
pair estimate holds on $\Fclass(N,P,s+1,\cdot,\epsilon)$. We will show that
we may take $P_1=P+1$ and $\epsilon_1=\epsilon/3$.

Assume that
\[
 H\leq\min\left(b-a,\frac{2\epsilon_1N}{s+P_1}\right)
 =\min\left(b-a,\frac{2\epsilon N}{3(s+P+1)}\right).
\]
In the notation of (2.3.4), we have
\[
 |S|^2\ll\frac{N^2}{H}+\frac NH\sum_{1\leq h\leq H}|S_1(h)|.
\]
By Lemma 3.7,
\[
 f_1\in\Fclass(N,P,s+1,shy,\epsilon),
\]
so the exponent pair $(k,l)$ may be applied to $S_1(h)$. Consequently,
\[
\begin{aligned}
 |S|^2
 &\ll H^{-1}N^2+H^{-1}N\sum_{1\leq h\leq H}
 \left\{(hLN^{-1})^kN^l+h^{-1}L^{-1}N\right\} \\
 &\ll H^{-1}N^2+H^kL^kN^{l-k+1}
      +H^{-1}L^{-1}N^2\log N.
\end{aligned}
\]
Since we are assuming that $L\geq\log N$, the first term dominates the
third. Applying Lemma 2.4 and using the upper bound on $H$ yields
\[
 |S|\ll L^\kappa N^\lambda+N(b-a)^{-1/2}.
\]
If the first term dominates, we are done. Otherwise, we employ the trivial
estimate to get
\[
 |S|\ll\min\{N(b-a)^{-1/2},b-a\}\ll N^{2/3}.
\]
Since $k\leq1/2$ and $l\geq1/2$, we have
\[
 \lambda=\frac12+\frac{l}{2k+2}\geq\frac23,
\]
and the desired estimate $|S|\ll L^\kappa N^\lambda$ follows.

The remaining case $1\leq L<\log N$ is easily dispatched. By Theorem 2.2,
\[
 |S|\ll N^{1/2}(\log N)^{1/2}\ll N^{2/3}
 \ll L^\kappa N^\lambda.
\]

\subsection*{3.5 PROOF OF THE B-PROCESS}

As in the previous section, we begin by showing that the needed auxiliary
function is in an appropriate $\Fclass$-set.

\medskip
\noindent\textbf{LEMMA 3.9.}
Suppose $P\geq2$, $f\in\Fclass(N,P,s,y,\epsilon)$, and $f$ is defined on the
interval $[a,b]$. Let $\alpha=f'(b)$ and $\beta=f'(a)$. For
$\nu\in[\alpha,\beta]$, let $x_\nu$ be defined by the relation
$f'(x_\nu)=\nu$, and define
\[
 \phi(\nu)=\nu x_\nu-f(x_\nu).
\]
Let $\sigma=1/s$ and $\eta=y^\sigma$. Then there is some constant
$C=C(s,P)$ such that
\[
 \left|\phi^{(p+1)}(\nu)-(-1)^p(\sigma)_p\eta\nu^{-\sigma-p}\right|
 <C\epsilon(\sigma)_p\eta\nu^{-\sigma-p}
\]
whenever $0\leq p\leq P-1$ and $\alpha\leq\nu\leq\beta$.

\noindent\textit{Proof.}
As in the proof of Lemma 3.7, let
\[
 F(x)=
 \begin{cases}
  yx^{1-s}(1-s)^{-1},&s\neq1,\\
  y\log x,&s=1.
 \end{cases}
\]
Condition (3.3.3) may be written as
\begin{equation*}
 |f^{(p+1)}(x)-F^{(p+1)}(x)|
 <\epsilon|F^{(p+1)}(x)|. \tag{3.5.1}
\end{equation*}
Now define $X_\nu$ by the relation $F'(X_\nu)=\nu$, and define
$\Phi(\nu)=\nu X_\nu-F(X_\nu)$. It is easy to see that
\[
 x_\nu=\phi'(\nu),
 \qquad
 X_\nu=\Phi'(\nu)=\eta\nu^{-\sigma}.
\]
To complete the proof, we will show that there is some constant $C=C(s,P)$
such that
\begin{equation*}
 |\phi^{(p+1)}(\nu)-\Phi^{(p+1)}(\nu)|
 <C\epsilon|\Phi^{(p+1)}(\nu)| \tag{3.5.2}
\end{equation*}
whenever $0\leq p\leq P-1$ and $\alpha\leq\nu\leq\beta$.

From our definitions, we note that $f'(\phi'(\nu))=\nu$. By differentiating
both sides with respect to $\nu$, we get
\[
 \phi''(\nu)=\frac1{f''(x_\nu)}.
\]
An easy induction argument shows that for $p\geq1$, we can find constants
$w(u_1,\ldots,u_{p-1})$ such that
\begin{equation*}
 \phi^{(p+1)}(\nu)
 =\frac{1}{\{f''(x_\nu)\}^{2p-1}}
  \sum_{\mathbf u}w(u_1,\ldots,u_{p-1})
  \prod_{j=1}^{p-1}f^{(u_j+1)}(x_\nu), \tag{3.5.3}
\end{equation*}
where the sum is over all $u_1,\ldots,u_{p-1}$ such that
$1\leq u_j\leq p$ and $u_1+\cdots+u_{p-1}=2p-2$. For $p=1$, the sum and
product are interpreted in the usual empty-product sense. The constants
$w(u_1,\ldots,u_{p-1})$ depend only on $u_1,\ldots,u_{p-1}$ and not on $f$,
so we also have
\begin{equation*}
 \Phi^{(p+1)}(\nu)
 =\frac{1}{\{F''(X_\nu)\}^{2p-1}}
  \sum_{\mathbf u}w(u_1,\ldots,u_{p-1})
  \prod_{j=1}^{p-1}F^{(u_j+1)}(X_\nu). \tag{3.5.4}
\end{equation*}

From (3.5.1), we see that
\[
 |f'(x_\nu)-F'(x_\nu)|
 =|\nu-yx_\nu^{-s}|<\epsilon yx_\nu^{-s}.
\]
Consequently,
\[
 (1-\epsilon)^\sigma\eta\nu^{-\sigma}
 <x_\nu
 <(1+\epsilon)^\sigma\eta\nu^{-\sigma},
\]
and so
\[
 |x_\nu-X_\nu|\ll\epsilon\eta\nu^{-\sigma}.
\]
Here, and in the remainder of this proof, all $\ll$-constants depend at most
on $s$ and $P$. From (3.5.1), we see that if $1\leq p\leq P-1$, then
\[
 |f^{(p+1)}(x_\nu)-F^{(p+1)}(x_\nu)|
 <\epsilon|F^{(p+1)}(x_\nu)|
 \ll\epsilon yN^{-s-p}.
\]
Furthermore, there is some $t$ in the open interval with endpoints $x_\nu$
and $X_\nu$ such that
\[
\begin{aligned}
 |F^{(p+1)}(x_\nu)-F^{(p+1)}(X_\nu)|
 &=|F^{(p+2)}(t)(X_\nu-x_\nu)| \\
 &\ll \epsilon yN^{-s-p-1}\eta\nu^{-\sigma}
 \ll \epsilon yN^{-s-p}.
\end{aligned}
\]
The last two estimates together yield
\[
 |f^{(p+1)}(x_\nu)-F^{(p+1)}(X_\nu)|
 \ll\epsilon yN^{-s-p}.
\]
Upon combining this with (3.5.3) and (3.5.4), we find, for $p\geq1$, that
\[
\begin{aligned}
 |\phi^{(p+1)}(\nu)-\Phi^{(p+1)}(\nu)|
 &\ll\frac{\epsilon y^{p-1}}{|F''(X_\nu)|^{2p-1}}
 \sum_{\mathbf u}|w(u_1,\ldots,u_{p-1})|\\
 &\qquad\qquad\times
 X_\nu^{-s(p-1)-(u_1+\cdots+u_{p-1})} \\
 &\ll\epsilon y^{-p}N^{ps+1}
 \ll\epsilon|\Phi^{(p+1)}(\nu)|.
\end{aligned}
\]
For $p=0$, (3.5.2) follows directly from the estimate for
$|x_\nu-X_\nu|$. This proves (3.5.2), and thus completes the proof of the
lemma.

In terms of our definition of $\Fclass$, Lemma 3.9 says that if
$f\in\Fclass(N,P,s,y,\epsilon)$, then the restriction of $\phi$ to the
interval $[\alpha,\beta]\cap[J,2J]$ is in
$\Fclass(J,P,\sigma,\eta,C\epsilon)$ for any $J$ with
$\alpha\leq J\leq\beta$.

\medskip
\noindent\textbf{THEOREM 3.10.}
If $(k,l)$ is an exponent pair, then
\[
 (\kappa,\lambda)=B(k,l)=(l-1/2,k+1/2)
\]
is an exponent pair.

\noindent\textit{Proof.}
First, $0\leq\kappa\leq1/2$ and $1/2\leq\lambda\leq1$ follow immediately
from $0\leq k\leq1/2\leq l\leq1$. Moreover, if $l=1/2$, then, by the remarks
at the end of Section 3.3, we have $k=1/2$. It follows that
$(\kappa,\lambda)=(0,1)$ is an exponent pair by the trivial estimate. We may
henceforth assume that $l>1/2$, and hence that $\kappa>0$.

Assume that $y>0$, $s>0$, $N>0$, and $L=yN^{-s}\geq1$. We want to find
$P_1$ and $\epsilon_1$ such that if
$f\in\Fclass(N,P_1,s,y,\epsilon_1)$, then
\[
 S=\sum_{n\in I}e(f(n))\ll L^\kappa N^\lambda.
\]
Put $\sigma=1/s$ and $\eta=y^\sigma$. Since $(k,l)$ is an exponent pair,
there exist $P>0$ and $\epsilon$ with $0<\epsilon<1/2$ such that the defining
estimate (3.3.4) holds for the parameter $\sigma$. By increasing $P$ if
necessary, we may assume that $P\geq4$. We take $P_1=P$ and
\[
 \epsilon_1=\min\left(\frac14,\frac{\epsilon}{C(s,P)}\right),
\]
where $C(s,P)$ is the constant occurring in Lemma 3.9.

Since $f$ satisfies the hypotheses of Lemma 3.6 with $F\asymp LN$, we may
write
\begin{equation*}
 S=\sum_{\alpha\leq\nu\leq\beta}
 \frac{e(-\phi(\nu)-1/8)}{|f''(x_\nu)|^{1/2}}
 +O\left(\log(2L)+L^{-1/2}N^{1/2}\right). \tag{3.5.5}
\end{equation*}
Here $\alpha=f'(b)$ and $\beta=f'(a)$, and
$\alpha\asymp\beta\asymp L$, with constants depending only on $s$.
Consequently, $[\alpha,\beta]$ can be covered by a bounded number of intervals
contained in dyadic intervals $[J,2J]$ with $J\asymp L$.

On each such interval, Lemma 3.9 and the exponent-pair hypothesis give,
uniformly in $w$,
\[
\begin{aligned}
 T_J(w)
 &\defeq\sum_{\substack{\nu\text{ in the interval}\\\nu\leq w}}
 e(\phi(\nu)) \\
 &\ll(\eta J^{-\sigma})^kJ^l+\eta^{-1}J^\sigma
 \ll N^kL^l+N^{-1}.
\end{aligned}
\]
Let
\[
 A(w)=|f''(x_w)|^{-1/2}.
\]
The defining estimates for $\Fclass$ imply
\[
 A(w)\asymp(LN^{-1})^{-1/2}
\]
and, since $x_w'=1/f''(x_w)$,
\[
 |A'(w)|\ll |f^{(3)}(x_w)|\,|f''(x_w)|^{-5/2}
 \ll N^{1/2}L^{-3/2}.
\]
Thus, on each of the above intervals,
\[
 \sup_w A(w)+\int|A'(w)|\,dw\ll(LN^{-1})^{-1/2}.
\]
Partial summation now gives
\[
\begin{aligned}
 \left|\sum_{\nu\text{ in the interval}}
 A(\nu)e(-\phi(\nu))\right|
 &=\left|\int A(w)\,d\overline{T_J(w)}\right| \\
 &\ll(N^kL^l+N^{-1})(LN^{-1})^{-1/2} \\
 &\ll L^\kappa N^\lambda+L^{-1/2}N^{-1/2}.
\end{aligned}
\]
Summing over the bounded number of intervals and using (3.5.5), we obtain
\[
 |S|\ll L^\kappa N^\lambda
       +\log(2L)+L^{-1/2}N^{1/2}.
\]
Since $\kappa>0$ and $L\geq1$, the first term dominates the remaining terms,
and the result is proved.

\subsection*{3.6 NOTES}

Here, we have presented the method of exponent pairs along very much the same
lines as Phillips (1933). One difference is in the use of Lemma 3.4, which is
somewhat stronger than the corresponding result of Phillips. Several variants
of this result have appeared in the literature; see, for example, Lemma 4.6
and pages 90--91 of Titchmarsh (1951), Lemma 10 of Heath-Brown (1983), and
pages 86--91 of Vinogradov (1985). The proof given here is due to H. L.
Montgomery (unpublished).

\end{document}
